% =============================================================
% Halaman Judul (title page), based on Lampiran 4 of the Pedoman.
% Identical to the cover, plus the "diajukan untuk..." purpose line.
% =============================================================
% See the note in cover.tex: the \titlepage environment is avoided
% because it resets \c@page, which must stay on the roman sequence here.
\clearpage
\thispagestyle{empty}
% See the note in cover.tex: 60mm here is measured from the paper edge.
\newgeometry{margin=0mm}
\begin{center}
\vspace*{60mm minus 10mm}

{\fontsize{14}{18}\selectfont\bfseries JUDUL TESIS UNTUK PROGRAM MAGISTER (S2)\\
DATA SCIENCE DAN KECERDASAN BUATAN\par}

\vspace{14mm}

{\fontsize{14}{18}\selectfont\bfseries TESIS\par}

\vspace{6mm}

{\fontsize{12}{16}\selectfont
Diajukan untuk melengkapi tugas dan memenuhi syarat memperoleh ijazah\\
Magister Data Science dan Kecerdasan Buatan\par}

\vspace{14mm}

{\fontsize{14}{18}\selectfont\bfseries NAMA MAHASISWA\\
NIM 000000000\par}

\vfill

\ThesisLogo

\vfill

{\fontsize{14}{18}\selectfont\bfseries
PROGRAM STUDI S2 DATA SCIENCE DAN KECERDASAN BUATAN\\
FAKULTAS ILMU KOMPUTER DAN TEKNOLOGI INFORMASI\\
UNIVERSITAS SUMATERA UTARA\\
MEDAN\\
\the\year\par}

\vspace*{60mm minus 10mm}
\end{center}
\restoregeometry
\clearpage
